\cotefile{53}
\scanpage{29}
% @ 1–4
d'un monoïde affine $G$, celà signifie aussi qu'il existe une opération de $G$ sur $M'$ \deleted{\unsure{qui}} telle que $M\unsure{-|\to}M'$ soit comp. avec \unsure{son|ses} op., et il est clair \inserted{\unsure{directemt}} (comme $W(M)\mathrel{-}W(M')$ \unsure{épi|épis}) que cette opération sera unique.

% @ 5–6
Si on a un mono
\[
M''\overset{i}{\hookrightarrow}M
\]
% @ 7–9
alors il ne sera pas vrai en général que $V(M)\mathrel{-}V(\unsure{M''|M'})$ soit surjectif, ni que $W(M'')\mathrel{-}W(M)$ soit injectif, donc \deleted{\unsure{faute de}} \unsure{cette|en tte} condition on \unsure{aura|a} du mal à parler d'opération induite.
% @ 10–12
\deleted{\unsure{Crt il est|Cert il est}} On dira qu'il \uline{existe} une op. induite si on peut faire opérer $G$ sur $M''$ de telle façon que $M''\mathrel{-}M$ soit comp. aux opérations.
% @ 13–15
Mais il n'est plus clair que l'op. ``induite'' soit unique ! En termes d'alg. comm., la question est de trouver un $\Delta_{M''}:M''\to M''\otimes A$ rendant commutatif
% @ 16–18 diagram
\diagram{\begin{tikzcd}[ampersand replacement=\&,row sep=large,column sep=large]
 M'' \arrow[r,dash,"{u_{M''}\,?}",to path={.. controls +(right:4mm) and +(left:4mm) .. (\tikztotarget) \tikztonodes -- ++(-.3pt,-1.4pt)}] \& M''\otimes A \arrow[d,dash,"{i\otimes{\mathrm{id}}_A}",to path={-- (\tikztotarget) \tikztonodes -- ++(.7pt,-1.2pt)}] \\
 M \arrow[u,dash,"i"',to path={-- (\tikztotarget) \tikztonodes .. controls +(up:2pt) and +(up:1pt) .. ++(2.5pt,.2pt)}] \arrow[r,dash,"{u_M}",to path={-- (\tikztotarget) \tikztonodes -- ++(0pt,1.1pt)}] \& M\otimes A
\end{tikzcd}}

% @ 19–20
On peut dire que c'est uniquement déterminé pourvu que $i\otimes{\mathrm{id}}_A=i_A$ soit un mono,
% @ 20 interlinear
\inserted{et si $\unsure{i|j}\otimes{\mathrm{id}}_{A\otimes A}$ l'est aussi, alors la condition $u_M(M'')\subset\mathrm{Im}\,M''_{\otimes A}$ est néc. et \uline{suffisante} pour la \deleted{\unsure{stabilité}} \unsure{stabilité|extensibilité} ; alors $M'$ est \unsure{stble|stbl} sss $M''$ l'est.}
% @ 21
\deleted{\unsure{batise|banise} \unsure{ces remarques}}. Les conditions envisagées sont OK dans les cas suivants :

% @ 22
a) $A$ est plat\unsure{$/k$|(le} (\unsure{p.~ex.} $k$ un \unsure{corps})

% @ 22–24 right margin
\margin{\deleted{\unsure{Alors|Mais} \unsure{les} \unsure{c.s.} de \unsure{la} \unsure{q°} \unsure{s'appl.} \unsure{montrent} $u_M(M'')\subset\mathrm{Im}\,M''\otimes A$.} Alors l'intersection \deleted{\unsure{finie}} des sous-modules stables est stable, \unsure{soit|et} \unsure{leur|l'inf} \unsure{somme|pour} \unsure{quelconques|quelconque}.}
% @ 23
b) $M/M''$ est \inserted{$A$} plat

% @ 24
c) $M''$ facteur direct dans $M$.

% @ 24 left margin
% slanted
\margin{\deleted{Supposons ces conditions vérifiées, alors $M'$ stable sss $M''$ stable}}
% @ 25
\deleted{\unsure{ou} $A$ \unsure{libre bien}.}

% @ 26–29
Dans le cas $A$ libre [plus gén. somme de modules proj. de t. f.] on voit que $M'$ est stable sss $M'$ est \deleted{stable sous l'action de $\check{A}$.} un sous-$\check{A}$-module. On trouve donc

% @ 30–32
\uline{Proposition} Si $A$ est somme de modules proj. de t. f., alors un sous-$k$-module \deleted{\unsure{de|du}} $A$-comodule est \unsure{stable} sss il est \deleted{stable sous l'op.} \unsure{un|est} sous $\unsure{\tilde{A}|\check{A}}$-module.
% archive sheet number: 28
